% GENERATED FILE. Edit canonical lessons, then run npm run generate:books.
% canonical-source-hash: fnv1a64:60585b7a2b8377bb
% canonical-lessons: MR-C529-samna, MR-C529-karnadhar, MR-C529-panc, MR-C529-maidan, MR-C529-mat-karne

\chapter{A Day at the Match}
\label{ch:mr-a-day-at-the-match}
\hlchaptermodality{529}

\begin{chapteropening}
\textbf{By the end of this chapter:} \emph{I can talk about a match, its captain and its umpire.}

Complete the last lesson of chapter 529: I can talk about a match, its captain and its umpire.
\end{chapteropening}

\section[\texorpdfstring{\mr{सामना} (sāmnā)}{सामना (sāmnā)}]{\mr{सामना} (sāmnā) — a match; a contest}

\label{lesson:MR-C529-samna}



\begin{quote}
Before the new one: say the Marathi for entertainment, then the Marathi for to be absorbed.
\end{quote}

\subsection*{You'll want to know: \mr{सामना}}
\textbf{\mr{सामना}} — \emph{sāmnā} — \textquotedblleft{}a match; a contest\textquotedblright{}.

\begin{quote}
\mr{काल} \mr{क्रिकेटचा} \mr{सामना} \mr{खूप} \mr{रंगला}.
\end{quote}

\begin{quote}
\emph{kāl krikeṭcā sāmnā khūp raṅglā.}
\end{quote}

\begin{quote}
Yesterday's cricket match was very exciting.
\end{quote}

Masculine.

The oblique is \mr{सामन्या}-: \mr{सामन्यात}.

\begin{cousinweb}
\textbf{\mr{सामना}} (\emph{sāmnā}) is usually linked with Hindi \emph{sāmnā}, \textquotedblleft{}a facing, a confrontation\textquotedblright{}. CDIAL sets the words for facing and in front under Sanskrit \emph{saṁmukha}, \textquotedblleft{}facing\textquotedblright{} (CDIAL 12982), but lists no \mr{सामना}, so the path is not settled.
\end{cousinweb}

\begin{grammarlens}[title={What you've built}]
One more word for describing what you lived and felt.
\end{grammarlens}

\subsection*{Your turn}
\begin{itemize}
\raggedright
  \item \emph{Say it:} \emph{sāmnā}
  \item \emph{Say it:} \emph{sāmnā}, once more
  \item \emph{Say it:} \emph{kāl krikeṭcā sāmnā khūp raṅglā.} — Yesterday's cricket match was very exciting.
  \item \emph{Recall:} say the Marathi for entertainment, then the Marathi for to be absorbed, then say \emph{sāmnā} again
\end{itemize}

\begin{tcolorbox}[breakable,colback=teal!4,colframe=teal!35!black,title={Before you move on}]
What does \mr{सामना} mean? (\textquotedblleft{}A match; a contest\textquotedblright{}.) Say it once more.
\end{tcolorbox}
\section[\texorpdfstring{\mr{कर्णधार} (karṇadhār)}{कर्णधार (karṇadhār)}]{\mr{कर्णधार} (karṇadhār) — a team captain}

\label{lesson:MR-C529-karnadhar}



\begin{quote}
Before the new one: say the Marathi for to be absorbed, then the Marathi for a match.
\end{quote}

\subsection*{You'll want to know: \mr{कर्णधार}}
\textbf{\mr{कर्णधार}} — \emph{karṇadhār} — \textquotedblleft{}a team captain\textquotedblright{}.

\begin{quote}
\mr{मीरा} \mr{आमच्या} \mr{संघाची} \mr{कर्णधार} \mr{आहे}.
\end{quote}

\begin{quote}
\emph{mīrā āmcyā saṅghācī karṇadhār āhe.}
\end{quote}

\begin{quote}
Meera is our team's captain.
\end{quote}

Masculine.

One word for a man or a woman.

\begin{cousinweb}
\textbf{\mr{कर्णधार}} (\emph{karṇadhār}) is a tatsama, taken from Sanskrit \emph{karṇadhāra}, \textquotedblleft{}helmsman\textquotedblright{} (Monier-Williams): the captain steers the team.
\end{cousinweb}

\begin{grammarlens}[title={What you've built}]
One more word for describing what you lived and felt.
\end{grammarlens}

\subsection*{Your turn}
\begin{itemize}
\raggedright
  \item \emph{Say it:} \emph{karṇadhār}
  \item \emph{Say it:} \emph{karṇadhār}, once more
  \item \emph{Say it:} \emph{mīrā āmcyā saṅghācī karṇadhār āhe.} — Meera is our team's captain.
  \item \emph{Recall:} say the Marathi for to be absorbed, then the Marathi for a match, then say \emph{karṇadhār} again
\end{itemize}

\begin{tcolorbox}[breakable,colback=teal!4,colframe=teal!35!black,title={Before you move on}]
What does \mr{कर्णधार} mean? (\textquotedblleft{}A team captain\textquotedblright{}.) Say it once more.
\end{tcolorbox}
\section[\texorpdfstring{\mr{पंच} (pañc)}{पंच (pañc)}]{\mr{पंच} (pañc) — an umpire; a referee}

\label{lesson:MR-C529-panc}



\begin{quote}
Before the new one: say the Marathi for a match, then the Marathi for a team captain.
\end{quote}

\subsection*{You'll want to know: \mr{पंच}}
\textbf{\mr{पंच}} — \emph{pañc} — \textquotedblleft{}an umpire; a referee\textquotedblright{}.

\begin{quote}
\mr{पंचांनी} \mr{पावसामुळे} \mr{खेळ} \mr{थांबवला}.
\end{quote}

\begin{quote}
\emph{pañcānnī pāvsāmuḷe kheḷ thāmbavlā.}
\end{quote}

\begin{quote}
The umpires stopped play because of the rain.
\end{quote}

Masculine.

A match has two, so often plural: \mr{पंचांनी}, the umpires.

\begin{cousinweb}
\textbf{\mr{पंच}} (\emph{pañc}) is a tadbhava of Sanskrit \emph{pañca}, \textquotedblleft{}five\textquotedblright{} (CDIAL 7655): a council of five once settled disputes, and the word came to mean a judge of fair play.
\end{cousinweb}

\begin{grammarlens}[title={What you've built}]
One more word for describing what you lived and felt.
\end{grammarlens}

\subsection*{Your turn}
\begin{itemize}
\raggedright
  \item \emph{Say it:} \emph{pañc}
  \item \emph{Say it:} \emph{pañc}, once more
  \item \emph{Say it:} \emph{pañcānnī pāvsāmuḷe kheḷ thāmbavlā.} — The umpires stopped play because of the rain.
  \item \emph{Recall:} say the Marathi for a match, then the Marathi for a team captain, then say \emph{pañc} again
\end{itemize}

\begin{tcolorbox}[breakable,colback=teal!4,colframe=teal!35!black,title={Before you move on}]
What does \mr{पंच} mean? (\textquotedblleft{}An umpire; a referee\textquotedblright{}.) Say it once more.
\end{tcolorbox}
\section[\texorpdfstring{\mr{मैदान} (maidān)}{मैदान (maidān)}]{\mr{मैदान} (maidān) — a playing field; a ground}

\label{lesson:MR-C529-maidan}



\begin{quote}
Before the new one: say the Marathi for a team captain, then the Marathi for an umpire.
\end{quote}

\subsection*{You'll want to know: \mr{मैदान}}
\textbf{\mr{मैदान}} — \emph{maidān} — \textquotedblleft{}a playing field; a ground\textquotedblright{}.

\begin{quote}
\mr{शाळेच्या} \mr{मैदानावर} \mr{रोज} \mr{संध्याकाळी} \mr{मुलं} \mr{खेळतात}.
\end{quote}

\begin{quote}
\emph{śāḷecyā maidānāvar roj sandhyākāḷī mulaṁ kheḷtāt.}
\end{quote}

\begin{quote}
Children play on the school ground every evening.
\end{quote}

Neuter.

\begin{cousinweb}
\textbf{\mr{मैदान}} (\emph{maidān}) is a loan from Persian \emph{maidān}, \textquotedblleft{}open ground, a field\textquotedblright{}.
\end{cousinweb}

\begin{grammarlens}[title={What you've built}]
One more word for describing what you lived and felt.
\end{grammarlens}

\subsection*{Your turn}
\begin{itemize}
\raggedright
  \item \emph{Say it:} \emph{maidān}
  \item \emph{Say it:} \emph{maidān}, once more
  \item \emph{Say it:} \emph{śāḷecyā maidānāvar roj sandhyākāḷī mulaṁ kheḷtāt.} — Children play on the school ground every evening.
  \item \emph{Recall:} say the Marathi for a team captain, then the Marathi for an umpire, then say \emph{maidān} again
\end{itemize}

\begin{tcolorbox}[breakable,colback=teal!4,colframe=teal!35!black,title={Before you move on}]
What does \mr{मैदान} mean? (\textquotedblleft{}A playing field; a ground\textquotedblright{}.) Say it once more.
\end{tcolorbox}
\section[\texorpdfstring{\mr{मात} \mr{करणे} (māt karṇe)}{मात करणे (māt karṇe)}]{\mr{मात} \mr{करणे} (māt karṇe) — to overcome; to beat}

\label{lesson:MR-C529-mat-karne}



\begin{quote}
Before the new one: say the Marathi for an umpire, then the Marathi for a playing field.
\end{quote}

\subsection*{You'll want to know: \mr{मात} \mr{करणे}}
\textbf{\mr{मात} \mr{करणे}} — \emph{māt karṇe} — \textquotedblleft{}to overcome; to beat\textquotedblright{}.

\begin{quote}
\mr{आमच्या} \mr{संघाने} \mr{मोठ्या} \mr{संघावर} \mr{मात} \mr{केली}.
\end{quote}

\begin{quote}
\emph{āmcyā saṅghāne moṭhyā saṅghāvar māt kelī.}
\end{quote}

\begin{quote}
Our team beat a bigger team.
\end{quote}

Past, with -\mr{ने} (he …): \mr{त्याने} \mr{अडचणींवर} \mr{मात} \mr{केली}, \emph{tyāne aḍcaṇīṁvar māt kelī}.

The one beaten takes -\mr{वर}: \mr{संघावर} \mr{मात} \mr{करणे}.

\begin{cousinweb}
\textbf{\mr{मात} \mr{करणे}} (\emph{māt karṇe}) is \mr{मात}, a loan from Persian \emph{māt}, \textquotedblleft{}checkmated\textquotedblright{}, usually traced to Arabic, with \mr{करणे}.
\end{cousinweb}

\begin{grammarlens}[title={What you've built}]
One more word for describing what you lived and felt.
\end{grammarlens}

\subsection*{Your turn}
\begin{itemize}
\raggedright
  \item \emph{Say it:} \emph{māt karṇe}
  \item \emph{Say it:} \emph{māt karṇe}, once more
  \item \emph{Say it:} \emph{āmcyā saṅghāne moṭhyā saṅghāvar māt kelī.} — Our team beat a bigger team.
  \item \emph{Recall:} say the Marathi for an umpire, then the Marathi for a playing field, then say \emph{māt karṇe} again
\end{itemize}

\begin{tcolorbox}[breakable,colback=teal!4,colframe=teal!35!black,title={Before you move on}]
What does \mr{मात} \mr{करणे} mean? (\textquotedblleft{}To overcome; to beat\textquotedblright{}.) Say it once more.
\end{tcolorbox}
